\documentclass[UTF8,11pt]{ctexart}

\usepackage[a4paper,margin=1in]{geometry}
\usepackage{amsmath}
\usepackage{array}
\usepackage{booktabs}
\usepackage{longtable}
\usepackage{enumitem}
\usepackage{xcolor}
\usepackage{hyperref}

\hypersetup{colorlinks=true,linkcolor=black,urlcolor=blue,citecolor=black}
\setlength{\parindent}{2em}
\setlength{\parskip}{0.25em}
\renewcommand{\arraystretch}{1.25}
\newcolumntype{L}[1]{>{\raggedright\arraybackslash}p{#1}}
\newcommand{\code}[1]{\texttt{\detokenize{#1}}}
\renewcommand{\abstractname}{摘要}
\emergencystretch=3em

\title{OpenPI 真机控制时序与同步假设分析}
\author{Codex}
\date{2026年5月31日}

\begin{document}
\maketitle

\begin{abstract}
本文分析目标仓库中 OpenPI 真机部署路径的推理、动作排程与底层执行机制，并将其与
《The Speedup Paradox: Rethinking Inference Optimization in Embodied Tasks》
中的静态同步假设进行对照。结论是：当前真机路径即使关闭显式异步推理，
也不是论文中理想化的严格同步执行，而是带有真实时间戳、过期动作丢弃、
动作截断和底层插值控制的时间戳排程式同步。因此，单纯放大推理延迟
\(T_{\mathrm{inf}}\) 会改变实际执行的动作索引、有效动作数和低层控制器接受的
waypoint，从而可能显著恶化控制质量。LIBERO 这类典型静态模拟器的默认评测
机制则更接近论文同步假设，因为仿真状态通常只在 \code{env.step} 时推进，
阻塞式推理不会被真实 wall-clock 时间转化为动作过期。
\end{abstract}

\section{问题背景}

论文的静态任务实验把推理延迟视作一个可以独立改变的变量。换言之，在观测相同、
模型相同、轻量化设置相同的条件下，如果仅改变 \(T_{\mathrm{inf}}\)，理论上改变的是
\(\rho = T_{\mathrm{inf}}/T_{\mathrm{act}}\)，而不是动作序列本身。这个假设适合推导
chunk 级时间收益，也适合在静态模拟器中隔离硬件速度的影响。

用户在 RTX 4070 上对 pick-and-place 任务进行去噪步数实验：10 步、5 步和 1 步
去噪使 \(T_{\mathrm{chunk}}\) 只有小幅下降，但动作质量随去噪步数减少而变差，
导致完成任务所需 chunk 数或动作数 \(N\) 明显上升，最终 \(T_{\mathrm{task}}\)
不降反升。按照论文理论，如果硬件算力变差、\(\rho\) 升高，则减少去噪步数带来的
\(T_{\mathrm{chunk}}\) 降幅应当更显著，可能出现某个中间步数的最优点，甚至最优点向
小步数方向移动。

实际放大 \(T_{\mathrm{inf}}\) 后，真机控制质量却严重下降，并出现过冲。这说明真实
部署中“改变推理延迟”并不只是改变 \(\rho\)。关键原因在于 OpenPI 真机路径把模型
输出映射到真实时间戳上的 future actions，并在推理完成后根据当前 wall-clock 时间
丢弃已经过期的动作。于是，延迟变大将改变被执行的 action index、有效执行数量
以及动作相对底层控制器的时间裕量。

\section{论文中的理想同步模型}

论文第 3 节给出的同步基线可写为
\begin{equation}
T_{\mathrm{chunk}} = T_{\mathrm{inf}} + nT_{\mathrm{act}},
\end{equation}
其中 \(T_{\mathrm{inf}}\) 是一次策略推理的耗时，\(n\) 是一个 chunk 中实际执行的
动作数，\(T_{\mathrm{act}}\) 是单个动作的执行时间。任务总时间写作
\begin{equation}
T_{\mathrm{task}} = N T_{\mathrm{chunk}},
\end{equation}
而硬件和控制频率的相对关系由
\begin{equation}
\rho = \frac{T_{\mathrm{inf}}}{T_{\mathrm{act}}}
\end{equation}
刻画。

这一模型隐含几个强假设。第一，推理阶段和执行阶段串行排列；推理完成后执行
固定数量的动作。第二，每轮实际执行的 \(n\) 不随 \(T_{\mathrm{inf}}\) 改变。
第三，模型输出动作的内容和被执行的动作索引不因推理延迟而改变。第四，对于静态
任务，推理期间环境状态不会发生与任务成功相关的变化。论文的 limitations 也明确
承认其分析抽象掉低层执行细节，并假设仅改变 inference latency 不会改变轨迹。

如果一个系统严格满足这些假设，那么放大 \(T_{\mathrm{inf}}\) 的作用主要体现在
\(T_{\mathrm{chunk}}\) 上；除非模型质量或轻量化程度改变，\(N\) 不应因为硬件延迟
本身发生系统性变化。当前真机框架的问题正在于：代码中存在多处把 \(T_{\mathrm{inf}}\)
转化为动作丢弃、动作重索引和底层排程裕量变化的机制。

\section{OpenPI 真机路径的执行语义}

\subsection{显式异步推理不是唯一的异步来源}

在 \code{deploy/inference_real.py} 中，默认参数包含
\code{frequency=10.0}、\code{steps_per_inference=6}、
\code{action_exec_latency=0.01}、\code{async_inference=False} 和
\code{inference_overlap_steps=0}。当 \code{async_inference=False} 时，主循环确实会
阻塞等待 \code{policy_infer(policy_obs)} 返回，因此不存在显式后台推理线程。
但是，这不等价于论文的严格同步。推理返回后，动作不是无条件从第 0 个开始执行，
而是先被映射到以观测时间为起点的时间戳序列，再按当前时间筛掉过期动作。

这种机制可抽象为
\begin{equation}
t_k = t_{\mathrm{obs}} + k\Delta t,\quad \Delta t = \frac{1}{f},
\end{equation}
并选择满足
\begin{equation}
t_k > t_{\mathrm{now}} + \delta_{\mathrm{exec}}
\end{equation}
的动作。若 \(T_{\mathrm{inf}}\) 增大，则 \(t_{\mathrm{now}}\) 后移，最小可执行索引
\(k_0\) 随之增大。于是模型输出的前若干动作会被直接跳过。实际执行的动作数近似为
\begin{equation}
n_{\mathrm{client}}
= \min\left(n_{\mathrm{cfg}}, H-k_0\right),
\end{equation}
其中 \(H\) 是模型 action horizon，\(n_{\mathrm{cfg}}\) 对应
\code{steps_per_inference}。因此，即使配置中的 \(n_{\mathrm{cfg}}\) 不变，有效执行
的首动作索引和动作集合也会随 \(T_{\mathrm{inf}}\) 改变。

\subsection{真实控制器还会再次处理时间戳}

动作通过 \code{env.exec_actions(..., compensate_latency=True)} 下发到 UMI 环境。
\code{bimanual_umi_env.py} 会再次根据接收时间过滤 future timestamps，并在补偿延迟
时从目标时间中减去 \code{robot_action_latency} 和 \code{gripper_action_latency}。
配置文件 \code{deploy/configs/umi_ur5e_wsg50.yaml} 中这两个延迟均为 0.1 秒，而
\code{inference_real.py} 默认的 \code{action_exec_latency} 只有 0.01 秒。若客户端
认为某动作仍在未来，但底层补偿后 waypoint 时间已经接近当前时间甚至落在过去，
RTDE 插值控制器可能无法按预期接收该 waypoint。

底层 UR 控制器不是在每个策略动作处停顿等待下一次推理，而是通过
\code{rtde_interpolation_controller.py} 中的 500Hz \code{servoL} 循环持续运行。
策略侧发送的是未来 waypoint；控制器侧将 wall-clock 时间转换为 monotonic 时间并
排入插值轨迹。\code{pose_trajectory_interpolator.py} 中对于时间不大于当前时刻的
waypoint 直接返回原轨迹，不会产生新的有效段。这进一步说明，放大
\(T_{\mathrm{inf}}\) 可能改变的不只是 chunk 时间，还包括底层轨迹是否真正被更新。

\subsection{延迟仿真会改变动作排程而非只改变等待时间}

\code{deploy/latency_simulation.py} 的 \code{scaled_infer} 先执行真实推理，再按照
耗时比例额外 sleep。也就是说，模型输出本身不变，但推理返回时的 wall-clock 时间
改变。对于论文模型，这相当于只改变 \(T_{\mathrm{inf}}\)。对于当前真机代码，这会
让 \code{_future_action_schedule} 看到更晚的 \code{time.time()}，从而丢弃更多前缀
动作，并改变实际执行的 action index。这个差异正好解释了为什么“控制模型轻量化程度
一定，只放大推理延迟”仍然会导致严重控制质量下降。

\section{代码证据}

表 \ref{tab:real-evidence} 汇总了支撑上述判断的关键代码位置。行号对应本次分析时
工作区中的文件内容。

\begin{longtable}{L{0.27\linewidth}L{0.36\linewidth}L{0.29\linewidth}}
\caption{OpenPI 真机路径中的时序机制证据}
\label{tab:real-evidence}\\
\toprule
代码位置 & 机制 & 对论文同步假设的影响 \\
\midrule
\endfirsthead
\toprule
代码位置 & 机制 & 对论文同步假设的影响 \\
\midrule
\endhead
\bottomrule
\endfoot
\code{deploy/inference_real.py:33--49} &
默认控制频率为 10Hz，\code{steps_per_inference} 默认为 6，
\code{action_exec_latency} 默认为 0.01 秒；显式异步推理默认关闭。 &
关闭 \code{async_inference} 只说明没有后台策略线程，不说明执行语义等于论文中的
严格同步。\\
\code{deploy/inference_real.py:160--178} &
\code{_future_action_schedule} 以 \code{obs_timestamp} 为起点给整个 action chunk
生成时间戳，并保留晚于 \code{time.time()+action_exec_latency} 的动作。 &
\(T_{\mathrm{inf}}\) 增大会让更多前缀动作被判定为过期，实际执行的动作索引随延迟
改变。\\
\code{deploy/inference_real.py:181--192} &
\code{_truncate_scheduled_actions} 将 future actions 截断到
\code{max_scheduled_actions} 或 \code{steps_per_inference}。 &
配置中的 \(n\) 是上限，不是实际保证执行的动作数量。\\
\code{deploy/inference_real.py:417--421} &
同步路径阻塞调用 \code{policy_infer(policy_obs)}，并记录 \code{infer_ms} 和
\code{obs_timestamp}。 &
策略推理是同步等待的，但推理完成后的动作选择仍依赖真实时间。\\
\code{deploy/inference_real.py:471--476} &
动作下发使用 \code{env.exec_actions(..., compensate_latency=True)}。 &
客户端选择的 future action 还会进入底层延迟补偿路径。\\
\code{deploy/inference_real.py:482--489} &
循环按配置中的计划步数推进 \code{step_idx}，而不是按底层实际接受 waypoint 数推进。 &
实验中观测到的有效 \(n\) 可能和循环索引中的 \(n\) 不一致。\\
\code{deploy/inference_real.py:491--518} &
异步模式通过 overlap steps 在动作执行期间提前发起下一轮推理。 &
若开启该模式，则系统进一步偏离论文的串行同步基线。\\
\code{deploy/latency_simulation.py:6--20} &
\code{scaled_infer} 在真实推理后按比例额外 sleep，并返回相同动作。 &
它看似只改变 \(T_{\mathrm{inf}}\)，但在真机代码中会改变 future action 过滤结果。\\
\code{deploy/async_inference.py:44--91} &
异步 worker 使用单个后台线程提交策略推理任务。 &
这是显式异步推理来源；不过即使关闭它，时间戳排程仍然存在。\\
\code{deploy/umi/real_world/bimanual_umi_env.py:420--431} &
\code{exec_actions} 按接收时间再次过滤 future timestamps。 &
客户端 future action 到达环境后仍可能被二次丢弃。\\
\code{deploy/umi/real_world/bimanual_umi_env.py:441--446} &
排程时从目标时间减去机器人和夹爪动作延迟。 &
若客户端 guard 过小，补偿后的 waypoint 可能落到过近或过去。\\
\code{deploy/configs/umi_ur5e_wsg50.yaml} &
\code{robot_action_latency=0.1} 且 \code{gripper_action_latency=0.1}。 &
底层补偿量显著大于客户端默认 \code{action_exec_latency=0.01}。\\
\code{deploy/umi/real_world/rtde_interpolation_controller.py:299--330} &
UR 侧以 500Hz \code{servoL} 循环持续执行插值轨迹。 &
底层执行是连续控制系统，不是每个 chunk 后停住等待推理的离散系统。\\
\code{deploy/umi/real_world/rtde_interpolation_controller.py:391--407} &
目标时间被转换为 monotonic 时间后排入插值器。 &
实际可执行性取决于 waypoint 到达控制器时是否仍有足够未来时间。\\
\code{deploy/umi/common/pose_trajectory_interpolator.py:122--126} &
若 waypoint 时间不晚于当前时间，则返回原轨迹。 &
过晚到达的动作不会形成新的有效轨迹段。\\
\code{src/openpi/training/config.py:1114--1118} &
UMI pi0.5 pick-place 配置的 \code{action_horizon} 为 30。 &
模型输出长度 \(H=30\)，但真机每轮通常只执行其中一小段。\\
\end{longtable}

\section{遥测结果与机制解释}

现有实验摘要位于
\code{data/umi_step_p3_trials_analysis/plots/tisd_chunk_metrics_trimmed_1_7_10_summary.csv}。
表 \ref{tab:telemetry} 显示，在 4070 上减少去噪步数并未显著降低
\(T_{\mathrm{chunk}}\)，但显著增加了 \(N\)，因此 \(T_{\mathrm{task}}\) 上升。

\begin{table}[h]
\centering
\caption{去噪步数实验的初步摘要}
\label{tab:telemetry}
\begin{tabular}{rrrr}
\toprule
去噪步数 & \(T_{\mathrm{chunk}}\) 均值/秒 & \(N\) 均值 & \(T_{\mathrm{task}}\) 估计/秒 \\
\midrule
1  & 0.7912 & 57.4 & 45.42 \\
7  & 0.7873 & 49.2 & 38.74 \\
10 & 0.7894 & 47.2 & 37.26 \\
\bottomrule
\end{tabular}
\end{table}

另外，遥测中常见 \code{model_action_count=30}，而
\code{scheduled_action_count} 约为 8，\code{dropped_action_count} 约为 22，
首个被排程动作的索引经常为 4 或 5。这与代码机制一致：模型生成了完整 horizon，
但客户端只挑出仍在未来的一段动作，再截断后下发。若人为放大
\(T_{\mathrm{inf}}\)，首个可执行索引会继续后移；在 pick-and-place 这类需要精细
接近和闭合的任务中，跳过前缀动作会造成动作相位错位、夹爪闭合时机偏移和末端过冲。

因此，在当前真机部署路径中，观测到的 \(N\) 上升不应只归因于去噪步数改变带来的
模型动作质量下降。额外的推理延迟还可能通过 action dropping 和 waypoint 过期机制
降低控制质量。换言之，真实系统中的 \(N\) 可以写作
\begin{equation}
N = N(q_{\mathrm{model}}, k_0, n_{\mathrm{client}}, n_{\mathrm{accepted}},
\epsilon_{\mathrm{schedule}}),
\end{equation}
其中 \(q_{\mathrm{model}}\) 表示模型动作质量，\(k_0\) 表示实际起始动作索引，
\(n_{\mathrm{client}}\) 表示客户端下发动作数，\(n_{\mathrm{accepted}}\) 表示底层控制器
实际接受的 waypoint 数，\(\epsilon_{\mathrm{schedule}}\) 表示排程误差。论文静态模型
主要把 \(N\) 看作模型质量和任务难度的结果，而当前真机路径使 \(N\) 同时依赖
wall-clock 时序。

\section{LIBERO 与论文假设的距离}

典型 LIBERO 评测机制更接近论文中的同步假设。以
\code{examples/libero/main.py} 为例，默认 \code{replan_steps=5}。程序只在
\code{action_plan} 为空时调用策略生成 action chunk，然后取
\code{action_chunk[:replan_steps]} 放入队列。每一轮从队列中弹出一个动作并调用
\code{env.step(action.tolist())}。在这个过程中，仿真状态只随 \code{env.step}
推进；阻塞等待策略推理不会让模拟世界按真实时间继续运动，也不会让旧动作因
wall-clock 时间过期。

因此，LIBERO 默认机制与论文假设的关系可以概括如下。它不是完全等同于论文，因为
实际 \(n\) 是 \code{replan_steps}，不一定等于模型的 action horizon；而且不同评测
脚本可能加入自己的缓存或并行化策略。但相对于真机 OpenPI，LIBERO 少了真实时间戳
过滤、底层控制器延迟补偿、网络传输、500Hz 插值控制器和 waypoint 过期逻辑。若在
LIBERO 中人为 sleep 以放大 \(T_{\mathrm{inf}}\)，通常只会增加每次 replan 的等待
时间，而不会自动跳过 action chunk 的前缀动作。这正是它更接近论文干净同步假设的
核心原因。

\section{为什么真机不默认采用严格同步}

真机和 LIBERO/论文假设拉开距离并不是任意设计，而是由真实机器人控制条件决定的。
第一，真实机械臂必须持续接收平滑轨迹或伺服目标。若策略推理期间机器人完全停等，
末端会频繁进入 stop-and-go 状态，导致速度不连续、接触不稳定，并可能触发底层控制
保护。第二，相机、网络、策略服务器和机器人控制器之间存在不可忽略且波动的延迟，
系统必须用未来时间戳和 latency compensation 把动作安排到控制器可执行的未来。
第三，真实世界在推理期间不会冻结；物体、夹爪、末端和传感器状态都在持续变化。
第四，底层 RTDE 控制器以高频连续插值执行，策略频率只有约 10Hz，两者天然需要一个
时序桥接层。

因此，真机系统更关心连续、平滑、可抢占和可容忍延迟抖动的控制，而论文静态同步
假设更关心隔离 \(T_{\mathrm{inf}}\)、\(T_{\mathrm{act}}\)、\(n\) 和 \(N\) 的理论关系。
严格同步可以作为受控实验模式，但不一定是生产式真机控制的最佳默认值。

\section{所谓同步推理平台中的异步层}

更一般地说，不能简单断言“所有号称同步推理的真机平台都在做异步推理”。更准确的
区分是：许多平台在策略 API 层确实是同步阻塞的，即用户代码调用 \code{infer()} 或
发送一次请求后等待返回；但在机器人系统层，它们几乎总会保留某种异步通信、轨迹缓冲、
时间戳排程或高频控制环。因此，“同步推理”通常只描述 policy-call 层面的阻塞语义，
不等价于论文中“推理期间物理世界冻结，推理返回后从 action chunk 第 0 个动作开始
严格执行固定 \(n\) 个动作”的理想同步。

这种差异在主流机器人中很常见。ROS 2 将 service 描述为短任务的同步 request/response
交互，而 action 则面向长时间任务，带 feedback、result 和 cancel/preempt 语义；
官方文档也明确指出 action server 在接收 goal 后执行任务并持续反馈，action 适用于
可能耗时较久的操作\footnote{\url{https://docs.ros.org/en/rolling/Concepts/Basic/About-Actions.html}}。
因此，即使 action client 选择阻塞等待 result，server 侧运动执行仍是一个独立的
长过程，而不是论文式的单线程停等。Universal Robots 的 RTDE 文档也把 RTDE 定义为
一种让外部应用通过 TCP/IP 与 UR controller 同步、同时不破坏控制器实时属性的接口，
并包含 setup 与 synchronization loop 两个阶段\footnote{\url{https://www.universal-robots.com/manuals/EN/HTML/SW10_10/Content/Prod-RTDE/Real_Time_Data_Exchange_RTDE.htm}}。
这说明外部策略程序和机器人控制器之间本来就是两个时钟域。Franka 的 libfranka 文档
同样说明 \code{franka::Robot::control} 会以 1kHz 频率执行用户回调，机器人状态也以
1kHz 提供\footnote{\url{https://support.franka.de/docs/libfranka.html}}。高层策略
是否阻塞等待推理，并不会消除这个底层实时控制循环。

因此，评估一个所谓“同步推理”的真机平台时，应该至少分清四层语义。第一是策略调用
是否同步：\code{infer()} 是否阻塞到模型返回。第二是动作下发是否同步：发送 trajectory
或 goal 后是否等待确认或结果。第三是控制器是否同步：底层是否仍在固定频率上插值、
伺服或跟踪上一段轨迹。第四才是论文式严格同步：推理期间不执行旧轨迹、不按真实时间
丢弃动作、不重排 action index，每轮都从第 0 个动作开始执行固定 \(n\) 个动作。主流
真机平台通常最多满足前两层，而第三层几乎必然是独立运行的；第四层更像受控实验模式，
需要显式实现和记录，而不是生产系统默认语义。

\section{当前实现：论文式严格同步与严格异步}

若目标是复现实验中“只改变 \(T_{\mathrm{inf}}\)，不改变动作索引和有效 \(n\)”的理想
同步或严格异步，应使用专门的实验模式，而不是复用当前 practical
\code{--async-inference} 路径。当前分支已经在 \code{deploy/inference_real.py}
中加入了两个互斥模式：\code{--strict-sync} 和 \code{--strict-async}。

\subsection{严格同步}

\code{--strict-sync} 的语义是：观察当前状态；阻塞推理；推理结束后从 action chunk
的第 0 个动作开始取固定数量 \(n\)；以推理完成后的当前时间为新的执行起点，将这
\(n\) 个动作重锚定到未来；等待这 \(n\) 个动作执行完，再进入下一轮观测和推理。
该模式绕过 \code{_future_action_schedule} 对 \code{obs_timestamp} 的依赖。时间戳为
\begin{equation}
t_k = t_{\mathrm{return}} + g + k\Delta t,\quad k=0,\ldots,n-1,
\end{equation}
其中 \(t_{\mathrm{return}}\) 是推理返回时间，\(g\) 对应
\code{--strict-sync-start-delay}，\(\Delta t=1/f\)。若仍启用
\code{compensate_latency=True}，则 \(g\) 必须至少覆盖 robot/gripper latency
和通信抖动；若希望最贴近论文假设，则可以在严格实验模式中关闭底层补偿，并选择一个
足够保守的固定 future start。关键点不是 guard 的具体数值，而是每次都执行同一组
动作索引 \(0,\ldots,n-1\)，且实际执行数量固定。

\subsection{严格异步}

\code{--strict-async} 对应论文中的固定 overlap 异步模型。设每个 chunk 固定执行
\(n\) 个动作，固定 overlap 步数为 \(n'\)，则下一轮推理在当前 chunk 执行到
\(n-n'\) 个动作后启动。若当前 chunk 的最后一个动作时间戳为 \(t_{n-1}\)，则下一
chunk 的计划起点为
\begin{equation}
t_{\mathrm{target}} = t_{n-1} + \Delta t,
\end{equation}
而下一轮推理启动时间为
\begin{equation}
t_{\mathrm{launch}} = t_{\mathrm{target}} - n'\Delta t.
\end{equation}
这样，模型推理和机械执行的理论 overlap 窗口固定为 \(n'\Delta t\)。若下一轮推理在
\(t_{\mathrm{target}}\) 前完成，则下一 chunk 贴着计划边界开始；若推理晚于边界完成，
则下一 chunk 起点顺延到
\begin{equation}
\max(t_{\mathrm{target}}, t_{\mathrm{return}} + g).
\end{equation}
在这两种情况下，下一 chunk 仍然从 action index 0 开始取固定数量动作，不会因为
wall-clock 过期而跳过 chunk 前缀。

该模式通过 \code{--strict-async-overlap-steps} 设置 \(n'\)，并要求
\(0 < n' \leq n\)。它与 \code{--async-inference} 和 \code{--strict-sync} 互斥，因为
三者代表不同实验语义：practical async 追求真机连续性，strict sync 追求完全串行
可控，strict async 追求固定 overlap 的论文式异步。

\subsection{遥测检查}

实验记录现在包含 \code{schedule_mode}、\code{scheduled_action_indices}、
\code{first_scheduled_action_index}、\code{last_scheduled_action_index} 和
\code{prefix_dropped_action_count}。严格异步还会记录
\code{strict_async_overlap_steps}、\code{strict_async_overlap_window_ms}、
\code{strict_async_target_start_timestamp}、\code{strict_async_launch_timestamp}
和 \code{strict_async_boundary_wait_ms}。只有当
\code{first_scheduled_action_index=0}、\code{prefix_dropped_action_count=0}
且固定动作数稳定时，该轮才真正满足论文式 strict sync/strict async 的动作索引条件。

\section{结论}

当前 OpenPI 真机框架的同步模式可以称为“阻塞推理 + 时间戳排程”，不能称为论文中的
严格同步。它的实际 \(\rho\) 不是简单的 \(T_{\mathrm{inf}}/T_{\mathrm{act}}\)，因为
\(T_{\mathrm{inf}}\) 同时影响 action dropping、动作起始索引、底层 waypoint 裕量和
有效 \(n\)。这解释了为什么在放大 \(T_{\mathrm{inf}}\) 时会出现控制质量严重下降甚至
过冲。

LIBERO 这类静态模拟器默认更接近论文假设，因为推理期间仿真状态通常不会按真实时间
推进，动作也不会因 wall-clock 延迟过期。若要在真机上验证论文理论，应使用当前新增的
\code{--strict-sync} 或 \code{--strict-async} 实验模式，把动作 chunk 在推理返回后
或固定 overlap 边界处重新锚定到未来，并保证每轮执行固定动作索引和固定数量。所谓
“同步推理”本身不足以保证这一点，因为它往往只约束策略调用层，而不约束底层控制器、
通信队列、trajectory buffer 和 timestamp scheduler。否则，实验测到的是模型质量、
硬件延迟、时间戳排程和底层控制器共同作用后的系统表现，而不是论文中单独改变
\(\rho\) 或固定 \(n'\) overlap 的理想化结果。

\end{document}
